% --- Archon physics formalization source begin ---
% archon:physics
% archon:covers PhyXMiniProblems/problem_phyx_mini_0754.lean
% archon:source-report reports/phyx_mini/problem_phyx_mini_0754.source.json

\chapter{Physics problem phyx\_mini\_0754}
\label{ch:PhyXMiniProblems_problem_phyx_mini_0754}

\paragraph{Problem source.}
\#\# Physical scenario

In figure, a ball is shot directly upward from the ground with an initial speed of \$v\_0 = 7.00\textbackslash{} m/s\$. Simultaneously, a construction elevator cab begins to move upward from the ground with a constant speed of \$v\_c = 3.00\textbackslash{} m/s\$.

\#\# Question

At what rate does the speed of the ball change relative to the cab floor?

\#\# Answer choices

- A: A: \$9.20\textbackslash{} m/s\textasciicircum{}2\$

- B: B: \$9.50\textbackslash{} m/s\textasciicircum{}2\$

- C: C: \$9.80\textbackslash{} m/s\textasciicircum{}2\$

- D: D: \$10.10\textbackslash{} m/s\textasciicircum{}2\$

\#\# Figure caption (auxiliary; use the image as primary evidence)

The image shows a man standing inside an elevator-like cage. The cage has a lattice pattern on its walls and is depicted as moving upwards, as indicated by an upward-pointing arrow labeled \textbackslash{}( \textbackslash{}vec\{v\}\_c \textbackslash{}).

To the right of the cage, there is a ball resting on a horizontal surface. An arrow labeled \textbackslash{}( \textbackslash{}vec\{v\}\_0 \textbackslash{}) points vertically upward from the ball, indicating its initial velocity. The ball is labeled clearly with the word "Ball."

The scene suggests a physics concept involving motion, likely focusing on the vertical velocities of both the cage and the ball.

\#\# Dataset metadata

Kinematics; Physical Model Grounding Reasoning, Spatial Relation Reasoning

\paragraph{Recorded answer/context.}
C: C: \$9.80\textbackslash{} m/s\textasciicircum{}2\$

\paragraph{Figure/image path.}
/root/proposal\_for\_physic/science-mango/phyx\_mini\_run/phyx\_data/test\_image/754.png

\paragraph{Formalization target.}
create a compiling Lean file with sorry bodies at `PhyXMiniProblems/problem\_phyx\_mini\_0754.lean`. The Lean declarations must preserve the physical quantities, dimensions or dimensional roles, figure labels, governing-law hypotheses, and final relation expressed by this problem.
Use Mathlib/Physlib names found through LeanExplore where available. If a physics API is missing, introduce faithful local abstractions rather than scalar placeholder aliases.

\begin{theorem}[Physics formalization target]
\label{thm:physics:phyx_mini_0754:target}
\lean{PhyXMiniProblems.ProblemPhyXMini0754.problem_phyx_mini_0754}
\uses{def:physics:phyx-mini-0754:phyxminiproblems-problemphyxmini0754-elevatorballsetup, def:physics:phyx-mini-0754:phyxminiproblems-problemphyxmini0754-matchesproblemstatement, def:physics:phyx-mini-0754:phyxminiproblems-problemphyxmini0754-matchesprimaryfigure, def:physics:phyx-mini-0754:phyxminiproblems-problemphyxmini0754-usesstandardterrestrialgravity, def:physics:phyx-mini-0754:phyxminiproblems-problemphyxmini0754-hasphysicalgravity, def:physics:phyx-mini-0754:phyxminiproblems-problemphyxmini0754-satisfieselevatorballkinematics, def:physics:phyx-mini-0754:phyxminiproblems-problemphyxmini0754-relativevelocitychangesatratemagnitudeat, def:physics:phyx-mini-0754:phyxminiproblems-problemphyxmini0754-answerchoiceaccelerationinmeterspersecondsquared}
The assigned autoformalize agent should translate this physics problem into Lean declarations in the covered file, with theorem and lemma proof bodies written as `by sorry`.
\end{theorem}
\begin{proof}
This is an autoformalization task, not a proof task. Produce statements that can later be proved by the physics prover without weakening the physical model.
\end{proof}

% --- Archon declaration topology begin ---
\section{Declaration topology}
\begin{definition}[Vertical Position Quantity]
  \label{def:physics:phyx-mini-0754:phyxminiproblems-problemphyxmini0754-verticalpositionquantity}
  \lean{PhyXMiniProblems.ProblemPhyXMini0754.VerticalPositionQuantity}
  A signed physical position along the vertical axis
\end{definition}
\begin{proof}
  This is the declared model definition of Vertical Position Quantity.
\end{proof}
\begin{definition}[Vertical Velocity Quantity]
  \label{def:physics:phyx-mini-0754:phyxminiproblems-problemphyxmini0754-verticalvelocityquantity}
  \lean{PhyXMiniProblems.ProblemPhyXMini0754.VerticalVelocityQuantity}
  A signed physical velocity along the vertical axis
\end{definition}
\begin{proof}
  This is the declared model definition of Vertical Velocity Quantity.
\end{proof}
\begin{definition}[Acceleration Magnitude Quantity]
  \label{def:physics:phyx-mini-0754:phyxminiproblems-problemphyxmini0754-accelerationmagnitudequantity}
  \lean{PhyXMiniProblems.ProblemPhyXMini0754.AccelerationMagnitudeQuantity}
  A physical acceleration magnitude
\end{definition}
\begin{proof}
  This is the declared model definition of Acceleration Magnitude Quantity.
\end{proof}
\begin{definition}[vertical Position Readout]
  \label{def:physics:phyx-mini-0754:phyxminiproblems-problemphyxmini0754-verticalpositionreadout}
  \lean{PhyXMiniProblems.ProblemPhyXMini0754.verticalPositionReadout}
  \uses{def:physics:phyx-mini-0754:phyxminiproblems-problemphyxmini0754-verticalpositionquantity}
  Read a signed position in the length unit selected by units
\end{definition}
\begin{proof}
  This is the declared model definition of vertical Position Readout.
\end{proof}
\begin{definition}[vertical Velocity Readout]
  \label{def:physics:phyx-mini-0754:phyxminiproblems-problemphyxmini0754-verticalvelocityreadout}
  \lean{PhyXMiniProblems.ProblemPhyXMini0754.verticalVelocityReadout}
  \uses{def:physics:phyx-mini-0754:phyxminiproblems-problemphyxmini0754-verticalvelocityquantity}
  Read a signed vertical velocity in the units selected by units
\end{definition}
\begin{proof}
  This is the declared model definition of vertical Velocity Readout.
\end{proof}
\begin{definition}[acceleration Magnitude Readout]
  \label{def:physics:phyx-mini-0754:phyxminiproblems-problemphyxmini0754-accelerationmagnitudereadout}
  \lean{PhyXMiniProblems.ProblemPhyXMini0754.accelerationMagnitudeReadout}
  \uses{def:physics:phyx-mini-0754:phyxminiproblems-problemphyxmini0754-accelerationmagnitudequantity}
  Read a nonnegative acceleration magnitude in the units selected by units
\end{definition}
\begin{proof}
  This is the declared model definition of acceleration Magnitude Readout.
\end{proof}
\begin{definition}[vertical Position In Meters]
  \label{def:physics:phyx-mini-0754:phyxminiproblems-problemphyxmini0754-verticalpositioninmeters}
  \lean{PhyXMiniProblems.ProblemPhyXMini0754.verticalPositionInMeters}
  \uses{def:physics:phyx-mini-0754:phyxminiproblems-problemphyxmini0754-verticalpositionquantity, def:physics:phyx-mini-0754:phyxminiproblems-problemphyxmini0754-verticalpositionreadout}
  Metre readout of a signed vertical position
\end{definition}
\begin{proof}
  This is the declared model definition of vertical Position In Meters.
\end{proof}
\begin{definition}[vertical Velocity In Meters Per Second]
  \label{def:physics:phyx-mini-0754:phyxminiproblems-problemphyxmini0754-verticalvelocityinmeterspersecond}
  \lean{PhyXMiniProblems.ProblemPhyXMini0754.verticalVelocityInMetersPerSecond}
  \uses{def:physics:phyx-mini-0754:phyxminiproblems-problemphyxmini0754-verticalvelocityquantity, def:physics:phyx-mini-0754:phyxminiproblems-problemphyxmini0754-verticalvelocityreadout}
  Metres-per-second readout of a signed vertical velocity
\end{definition}
\begin{proof}
  This is the declared model definition of vertical Velocity In Meters Per Second.
\end{proof}
\begin{definition}[acceleration Magnitude In Meters Per Second Squared]
  \label{def:physics:phyx-mini-0754:phyxminiproblems-problemphyxmini0754-accelerationmagnitudeinmeterspersecondsquared}
  \lean{PhyXMiniProblems.ProblemPhyXMini0754.accelerationMagnitudeInMetersPerSecondSquared}
  \uses{def:physics:phyx-mini-0754:phyxminiproblems-problemphyxmini0754-accelerationmagnitudequantity, def:physics:phyx-mini-0754:phyxminiproblems-problemphyxmini0754-accelerationmagnitudereadout}
  Metres-per-second-squared readout of an acceleration magnitude
\end{definition}
\begin{proof}
  This is the declared model definition of acceleration Magnitude In Meters Per Second Squared.
\end{proof}
\begin{definition}[Figure Velocity Arrow]
  \label{def:physics:phyx-mini-0754:phyxminiproblems-problemphyxmini0754-figurevelocityarrow}
  \lean{PhyXMiniProblems.ProblemPhyXMini0754.FigureVelocityArrow}
  The two velocity arrows printed in the supplied bitmap
\end{definition}
\begin{proof}
  This is the declared model definition of Figure Velocity Arrow.
\end{proof}
\begin{definition}[Figure Object]
  \label{def:physics:phyx-mini-0754:phyxminiproblems-problemphyxmini0754-figureobject}
  \lean{PhyXMiniProblems.ProblemPhyXMini0754.FigureObject}
  The objects whose initial locations are visible in the bitmap
\end{definition}
\begin{proof}
  This is the declared model definition of Figure Object.
\end{proof}
\begin{definition}[Vertical Direction]
  \label{def:physics:phyx-mini-0754:phyxminiproblems-problemphyxmini0754-verticaldirection}
  \lean{PhyXMiniProblems.ProblemPhyXMini0754.VerticalDirection}
  Orientation along the vertical axis
\end{definition}
\begin{proof}
  This is the declared model definition of Vertical Direction.
\end{proof}
\begin{definition}[Elevator Ball Figure]
  \label{def:physics:phyx-mini-0754:phyxminiproblems-problemphyxmini0754-elevatorballfigure}
  \lean{PhyXMiniProblems.ProblemPhyXMini0754.ElevatorBallFigure}
  \uses{def:physics:phyx-mini-0754:phyxminiproblems-problemphyxmini0754-figurevelocityarrow, def:physics:phyx-mini-0754:phyxminiproblems-problemphyxmini0754-figureobject, def:physics:phyx-mini-0754:phyxminiproblems-problemphyxmini0754-verticaldirection}
  Structured qualitative information read directly from the primary image
\end{definition}
\begin{proof}
  This is the declared model definition of Elevator Ball Figure.
\end{proof}
\begin{definition}[Elevator Ball Setup]
  \label{def:physics:phyx-mini-0754:phyxminiproblems-problemphyxmini0754-elevatorballsetup}
  \lean{PhyXMiniProblems.ProblemPhyXMini0754.ElevatorBallSetup}
  \uses{def:physics:phyx-mini-0754:phyxminiproblems-problemphyxmini0754-verticalpositionquantity, def:physics:phyx-mini-0754:phyxminiproblems-problemphyxmini0754-verticalvelocityquantity, def:physics:phyx-mini-0754:phyxminiproblems-problemphyxmini0754-accelerationmagnitudequantity, def:physics:phyx-mini-0754:phyxminiproblems-problemphyxmini0754-elevatorballfigure}
  The trajectories and dimensionful parameters of the one-dimensional model. The argument of each trajectory is an SI time readout in seconds. The relative-velocity trajectory is an independent physical quantity whose governing relation to the ball and cab velocities is stated below; no field assigns its derivative or the requested answer
\end{definition}
\begin{proof}
  This is the declared model definition of Elevator Ball Setup.
\end{proof}
\begin{definition}[Matches Problem Statement]
  \label{def:physics:phyx-mini-0754:phyxminiproblems-problemphyxmini0754-matchesproblemstatement}
  \lean{PhyXMiniProblems.ProblemPhyXMini0754.MatchesProblemStatement}
  \uses{def:physics:phyx-mini-0754:phyxminiproblems-problemphyxmini0754-verticalpositioninmeters, def:physics:phyx-mini-0754:phyxminiproblems-problemphyxmini0754-verticalvelocityinmeterspersecond, def:physics:phyx-mini-0754:phyxminiproblems-problemphyxmini0754-elevatorballsetup}
  Numerical values and simultaneous ground-level start stated in the problem. The initial velocity values are retained even though they cancel out of the relative acceleration
\end{definition}
\begin{proof}
  This is the declared model definition of Matches Problem Statement.
\end{proof}
\begin{definition}[Matches Primary Figure]
  \label{def:physics:phyx-mini-0754:phyxminiproblems-problemphyxmini0754-matchesprimaryfigure}
  \lean{PhyXMiniProblems.ProblemPhyXMini0754.MatchesPrimaryFigure}
  \uses{def:physics:phyx-mini-0754:phyxminiproblems-problemphyxmini0754-elevatorballsetup}
  Readouts from the supplied bitmap: both velocity arrows point upward, both objects begin on the ground line, the ball is labelled, and a person is drawn inside the construction cab
\end{definition}
\begin{proof}
  This is the declared model definition of Matches Primary Figure.
\end{proof}
\begin{definition}[Uses Standard Terrestrial Gravity]
  \label{def:physics:phyx-mini-0754:phyxminiproblems-problemphyxmini0754-usesstandardterrestrialgravity}
  \lean{PhyXMiniProblems.ProblemPhyXMini0754.UsesStandardTerrestrialGravity}
  \uses{def:physics:phyx-mini-0754:phyxminiproblems-problemphyxmini0754-accelerationmagnitudeinmeterspersecondsquared, def:physics:phyx-mini-0754:phyxminiproblems-problemphyxmini0754-elevatorballsetup}
  The conventional near-Earth gravitational calibration used by the recorded multiple-choice answer
\end{definition}
\begin{proof}
  This is the declared model definition of Uses Standard Terrestrial Gravity.
\end{proof}
\begin{definition}[Has Physical Gravity]
  \label{def:physics:phyx-mini-0754:phyxminiproblems-problemphyxmini0754-hasphysicalgravity}
  \lean{PhyXMiniProblems.ProblemPhyXMini0754.HasPhysicalGravity}
  \uses{def:physics:phyx-mini-0754:phyxminiproblems-problemphyxmini0754-accelerationmagnitudeinmeterspersecondsquared, def:physics:phyx-mini-0754:phyxminiproblems-problemphyxmini0754-elevatorballsetup}
  Positivity of the sole acceleration magnitude in the idealized model
\end{definition}
\begin{proof}
  This is the declared model definition of Has Physical Gravity.
\end{proof}
\begin{definition}[Has Vertical Velocity At]
  \label{def:physics:phyx-mini-0754:phyxminiproblems-problemphyxmini0754-hasverticalvelocityat}
  \lean{PhyXMiniProblems.ProblemPhyXMini0754.HasVerticalVelocityAt}
  \uses{def:physics:phyx-mini-0754:phyxminiproblems-problemphyxmini0754-verticalpositionquantity, def:physics:phyx-mini-0754:phyxminiproblems-problemphyxmini0754-verticalvelocityquantity, def:physics:phyx-mini-0754:phyxminiproblems-problemphyxmini0754-verticalpositioninmeters, def:physics:phyx-mini-0754:phyxminiproblems-problemphyxmini0754-verticalvelocityinmeterspersecond}
  At an SI time readout, the derivative of a metre-valued height readout is the corresponding metres-per-second velocity readout
\end{definition}
\begin{proof}
  This is the declared model definition of Has Vertical Velocity At.
\end{proof}
\begin{definition}[Has Vertical Acceleration Readout At]
  \label{def:physics:phyx-mini-0754:phyxminiproblems-problemphyxmini0754-hasverticalaccelerationreadoutat}
  \lean{PhyXMiniProblems.ProblemPhyXMini0754.HasVerticalAccelerationReadoutAt}
  \uses{def:physics:phyx-mini-0754:phyxminiproblems-problemphyxmini0754-verticalvelocityquantity, def:physics:phyx-mini-0754:phyxminiproblems-problemphyxmini0754-verticalvelocityinmeterspersecond}
  At an SI time readout, the derivative of a metres-per-second velocity readout has the given signed metres-per-second-squared value
\end{definition}
\begin{proof}
  This is the declared model definition of Has Vertical Acceleration Readout At.
\end{proof}
\begin{definition}[Satisfies Elevator Ball Kinematics]
  \label{def:physics:phyx-mini-0754:phyxminiproblems-problemphyxmini0754-satisfieselevatorballkinematics}
  \lean{PhyXMiniProblems.ProblemPhyXMini0754.SatisfiesElevatorBallKinematics}
  \uses{def:physics:phyx-mini-0754:phyxminiproblems-problemphyxmini0754-verticalvelocityreadout, def:physics:phyx-mini-0754:phyxminiproblems-problemphyxmini0754-accelerationmagnitudeinmeterspersecondsquared, def:physics:phyx-mini-0754:phyxminiproblems-problemphyxmini0754-elevatorballsetup, def:physics:phyx-mini-0754:phyxminiproblems-problemphyxmini0754-hasverticalvelocityat, def:physics:phyx-mini-0754:phyxminiproblems-problemphyxmini0754-hasverticalaccelerationreadoutat}
  Governing laws for ideal free fall beside a constant-speed cab. The ball accelerates downward with magnitude g; the cab velocity is constant; and relative velocity is ball velocity minus cab-floor velocity in every coherent unit system. None of these fields gives the magnitude of the relative acceleration or selects an answer choice
\end{definition}
\begin{proof}
  This is the declared model definition of Satisfies Elevator Ball Kinematics.
\end{proof}
\begin{lemma}[cab vertical velocity has zero derivative]
  \label{lem:physics:phyx-mini-0754:phyxminiproblems-problemphyxmini0754-cab-vertical-velocity-has-zero-derivative}
  \lean{PhyXMiniProblems.ProblemPhyXMini0754.cab_vertical_velocity_has_zero_derivative}
  \uses{def:physics:phyx-mini-0754:phyxminiproblems-problemphyxmini0754-elevatorballsetup, def:physics:phyx-mini-0754:phyxminiproblems-problemphyxmini0754-hasverticalaccelerationreadoutat, def:physics:phyx-mini-0754:phyxminiproblems-problemphyxmini0754-satisfieselevatorballkinematics}
  A constant cab velocity has zero signed acceleration readout
\end{lemma}
\begin{proof}
  The conclusion follows from the governing relations and figure readouts named in the statement of cab vertical velocity has zero derivative.
\end{proof}
\begin{lemma}[relative vertical velocity derivative]
  \label{lem:physics:phyx-mini-0754:phyxminiproblems-problemphyxmini0754-relative-vertical-velocity-derivative}
  \lean{PhyXMiniProblems.ProblemPhyXMini0754.relative_vertical_velocity_derivative}
  \uses{def:physics:phyx-mini-0754:phyxminiproblems-problemphyxmini0754-accelerationmagnitudeinmeterspersecondsquared, def:physics:phyx-mini-0754:phyxminiproblems-problemphyxmini0754-elevatorballsetup, def:physics:phyx-mini-0754:phyxminiproblems-problemphyxmini0754-hasverticalaccelerationreadoutat, def:physics:phyx-mini-0754:phyxminiproblems-problemphyxmini0754-satisfieselevatorballkinematics}
  Subtracting the zero cab acceleration from the ball's downward free-fall acceleration gives the signed derivative of the relative velocity
\end{lemma}
\begin{proof}
  The conclusion follows from the governing relations and figure readouts named in the statement of relative vertical velocity derivative.
\end{proof}
\begin{definition}[Relative Velocity Changes At Rate Magnitude At]
  \label{def:physics:phyx-mini-0754:phyxminiproblems-problemphyxmini0754-relativevelocitychangesatratemagnitudeat}
  \lean{PhyXMiniProblems.ProblemPhyXMini0754.RelativeVelocityChangesAtRateMagnitudeAt}
  \uses{def:physics:phyx-mini-0754:phyxminiproblems-problemphyxmini0754-elevatorballsetup, def:physics:phyx-mini-0754:phyxminiproblems-problemphyxmini0754-hasverticalaccelerationreadoutat}
  The magnitude of the signed derivative of the ball's velocity relative to the cab floor. This is the relative-acceleration magnitude meant by the source's phrase "rate does the speed ... change relative to the cab floor."
\end{definition}
\begin{proof}
  This is the declared model definition of Relative Velocity Changes At Rate Magnitude At.
\end{proof}
\begin{definition}[Answer Choice]
  \label{def:physics:phyx-mini-0754:phyxminiproblems-problemphyxmini0754-answerchoice}
  \lean{PhyXMiniProblems.ProblemPhyXMini0754.AnswerChoice}
  Labels of the four displayed multiple-choice answers
\end{definition}
\begin{proof}
  This is the declared model definition of Answer Choice.
\end{proof}
\begin{definition}[answer Choice Acceleration In Meters Per Second Squared]
  \label{def:physics:phyx-mini-0754:phyxminiproblems-problemphyxmini0754-answerchoiceaccelerationinmeterspersecondsquared}
  \lean{PhyXMiniProblems.ProblemPhyXMini0754.answerChoiceAccelerationInMetersPerSecondSquared}
  \uses{def:physics:phyx-mini-0754:phyxminiproblems-problemphyxmini0754-answerchoice}
  Exact SI acceleration readouts represented by the printed decimals
\end{definition}
\begin{proof}
  This is the declared model definition of answer Choice Acceleration In Meters Per Second Squared.
\end{proof}
\begin{definition}[recorded Answer Choice]
  \label{def:physics:phyx-mini-0754:phyxminiproblems-problemphyxmini0754-recordedanswerchoice}
  \lean{PhyXMiniProblems.ProblemPhyXMini0754.recordedAnswerChoice}
  \uses{def:physics:phyx-mini-0754:phyxminiproblems-problemphyxmini0754-answerchoice}
  The answer label recorded in the source dataset
\end{definition}
\begin{proof}
  This is the declared model definition of recorded Answer Choice.
\end{proof}
% --- Archon declaration topology end ---
% --- Archon physics formalization source end ---
